\section{Planar dynamic model}\label{sec:model}

The equations of motion of the robot in the vertical plane are derived with
the Euler--Lagrange method: the kinetic and potential energy are written as
functions of the generalised coordinates, the motor torque is introduced as a
generalised force, and the Euler--Lagrange equations give one equation of
motion per coordinate.

\subsection{Assumptions}

\begin{itemize}
  \item The robot moves in a straight line on flat, horizontal ground; the two
        wheels move together and are modelled as a single wheel.
  \item The robot consists of two rigid bodies: the body and the wheel.
  \item The wheel rolls without slipping.
  \item Friction, motor dynamics and tyre compliance are neglected.
  \item The motor is fixed to the body. It applies a torque $\tau$ to the
        wheel and, by action and reaction, a torque $-\tau$ to the body.
\end{itemize}

\begin{table}[h]
\centering
\caption{Parameters of the model.}
\label{tab:params}
\begin{tabular}{@{}lll@{}}
\toprule
Symbol & Description & Unit \\
\midrule
$\mb$ & mass of the body & \si{kg} \\
$\Ib$ & moment of inertia of the body about its centre of mass & \si{kg.m^2} \\
$l$   & distance between the wheel axle and the body's centre of mass & \si{m} \\
$\mw$ & mass of the wheel & \si{kg} \\
$\Iw$ & moment of inertia of the wheel about the axle & \si{kg.m^2} \\
$r$   & wheel radius & \si{m} \\
$g$   & gravitational acceleration & \si{m/s^2} \\
$\tau$ & motor torque applied to the wheel & \si{N.m} \\
\bottomrule
\end{tabular}
\end{table}

\subsection{Coordinates and conventions}

The inertial frame (\cref{fig:model}) has the $x$ axis horizontal, the $z$
axis vertical and pointing up, and the origin $O$ at the height of the wheel
axle; $y$ completes the right-handed frame and points into the page. The
configuration is described by:
\begin{itemize}
  \item $x$, the horizontal position of the wheel axle;
  \item $\theta$, the angle of the body from the vertical, positive when the
        top of the body tilts towards $+x$;
  \item $\varphi$, the rotation angle of the wheel with respect to the ground,
        positive when the wheel rolls towards $+x$.
\end{itemize}
Both angles are measured in radians and are positive clockwise in
\cref{fig:model}. In the drawing $\theta>0$ and $\varphi>0$.

\begin{figure}[h]
\centering
\begin{subfigure}[b]{\linewidth}
  \centering% Planar model, side view (robot seen from its right side)
\begin{tikzpicture}[scale=1.0]
  \def\th{18}   % drawn tilt angle, positive = top towards +x
  \def\R{1.0}   % wheel radius
  \def\L{2.0}   % axle -> body CoM
  \def\X{4.2}   % drawn axle position

  % ground
  \fill[pattern=north east lines,pattern color=gray] (-1.0,-\R) rectangle (6.4,-\R-0.2);
  \draw (-1.0,-\R) -- (6.4,-\R);

  % world frame at axle height
  \draw[->] (0,0) -- (1.2,0) node[below]{$x$};
  \draw[->] (0,0) -- (0,1.2) node[left]{$z$};
  \node[below left] at (0,0) {$O$};
  \draw (-0.55,0.35) circle (0.12);
  \draw (-0.63,0.27) -- (-0.47,0.43) (-0.63,0.43) -- (-0.47,0.27);
  \node[left] at (-0.68,0.35) {$y$};

  % position of the axle
  \draw (0,-\R-0.3) -- (0,-\R-0.7) (\X,-\R-0.3) -- (\X,-\R-0.7);
  \draw[<->] (0,-\R-0.5) -- node[below]{$x$} (\X,-\R-0.5);

  % body
  \begin{scope}[shift={(\X,0)},rotate=-\th]
    \draw[fill=gray!10] (-0.9,-0.35) rectangle (0.9,3.5);
    \draw[dashed,gray] (0,\L) -- (0,3.9);
  \end{scope}
  \draw[dashed,gray] (\X,0) -- (\X,3.9);
  \draw[->] (\X,3.25) arc[start angle=90,end angle=90-\th,radius=3.25];
  \node at ($(\X,0)+(90-\th/2:3.55)$) {$\theta$};

  % wheel
  \draw[fill=white] (\X,0) circle (\R);
  \fill (\X,0) circle (1.5pt);
  \draw (\X,0) -- ++(-125:\R);
  \draw[dashed,gray] (\X,0) -- ++(-90:\R);
  \draw[->] ($(\X,0)+(-90:0.6)$) arc[start angle=-90,end angle=-125,radius=0.6];
  \node at ($(\X,0)+(-150:0.62)$) {$\varphi$};
  \draw[<->] (\X,0) -- node[above=-1pt]{$r$} ++(165:\R);

  % axle -> CoM
  \begin{scope}[shift={(\X,0)},rotate=-\th]
    \draw[thick] (0,0) -- (0,\L) node[pos=0.75,left]{$l$};
    \fill (0,\L) circle (2.2pt);
    \node[left=3pt,fill=gray!10,inner sep=1pt] at (0,\L+0.45) {$\pb$};
  \end{scope}
  \draw[->] ($(\X,0)+(90-\th:\L)$) -- ++(0,-1.0) node[right,fill=white,inner sep=1pt]{$\mb g$};
\end{tikzpicture}

  \caption{Coordinates and parameters.}
  \label{fig:model-kin}
\end{subfigure}\\[1em]
\begin{subfigure}[b]{\linewidth}
  \centering% Motor torque pair and rolling without slipping
\begin{tikzpicture}[scale=1.0]
  \def\th{12}
  \def\R{1.0}
  \fill[pattern=north east lines,pattern color=gray] (-2.2,-\R) rectangle (4.0,-\R-0.2);
  \draw (-2.2,-\R) -- (4.0,-\R);

  \begin{scope}[rotate=-\th]
    \draw[fill=gray!10] (-0.9,-0.35) rectangle (0.9,2.6);
  \end{scope}
  \draw[->,thick] (60:1.75) arc[start angle=60,end angle=125,radius=1.75];
  \node[anchor=south,fill=white,inner sep=1pt] at (95:1.85) {$-\tau$ (on body)};

  \draw[fill=white] (0,0) circle (\R);
  \fill (0,0) circle (1.5pt);
  \draw[->,thick] (150:0.6) arc[start angle=150,end angle=20,radius=0.6];
  \node at (0,0.27) {$+\tau$};

  \draw[dashed,gray] (2.4,0) circle (\R);
  \draw[decorate,decoration={brace,amplitude=5pt,mirror}]
       (0,-\R-0.3) -- (2.4,-\R-0.3) node[midway,below=6pt]{$x = r\varphi$};
\end{tikzpicture}

  \caption{Motor torque pair and rolling without slipping.}
  \label{fig:model-act}
\end{subfigure}
\caption{Planar model of the robot.}
\label{fig:model}
\end{figure}

\subsection{Rolling constraint}

If the wheel does not slip, the distance travelled by the axle equals the arc
length unrolled on the ground (\cref{fig:model-act}). Choosing $\varphi=0$
when $x=0$,
\begin{equation}
  x = r\varphi .
  \label{eq:rolling}
\end{equation}
Differentiating with respect to time with the product rule, and using that the
radius is constant ($\dot r = 0$):
\begin{equation}
  \dot x = \frac{\dd (r\varphi)}{\dd t} = \dot r\,\varphi + r\dot\varphi = r\dot\varphi
  \qquad\Longrightarrow\qquad
  \dot\varphi = \frac{\dot x}{r}.
  \label{eq:rolling-vel}
\end{equation}
The constraint expresses $\varphi$ as a function of $x$, so the system has two
degrees of freedom. The generalised coordinates are
\begin{equation}
  \q = \begin{bmatrix} x \\ \theta \end{bmatrix},
  \qquad
  \qd = \begin{bmatrix} \dot x \\ \dot\theta \end{bmatrix}.
\end{equation}

\subsection{Position and velocity of the body's centre of mass}

The kinetic energy requires velocities measured in the inertial frame, so the
position of the body's centre of mass is written in that frame. It is the
position of the axle, $(x,0)$, plus a segment of length $l$ tilted by $\theta$
from the vertical, whose components are $l\sin\theta$ along $x$ and
$l\cos\theta$ along $z$:
\begin{equation}
  \pb = \begin{bmatrix} x + l\sin\theta \\ l\cos\theta \end{bmatrix}.
  \label{eq:pb}
\end{equation}
The velocity is the time derivative of each component. Both $x$ and $\theta$
are functions of time, while $l$ is constant:
\begin{equation}
  \vb = \dot{\bm p}_{\mathrm b}
  = \begin{bmatrix}
      \dfrac{\dd x}{\dd t} + \dfrac{\dd (l\sin\theta)}{\dd t} \\[10pt]
      \dfrac{\dd (l\cos\theta)}{\dd t}
    \end{bmatrix}.
\end{equation}
With the product rule, the term $\dot l\,\sin\theta$ vanishes because
$\dot l=0$; the chain rule then gives the derivative of a function of $\theta$
as its derivative with respect to $\theta$ times $\dot\theta$:
\begin{align}
  \frac{\dd (l\sin\theta)}{\dd t} &= \dot l\sin\theta + l\,\frac{\dd\sin\theta}{\dd\theta}\,\frac{\dd\theta}{\dd t} = l\cos\theta\,\dot\theta , \\
  \frac{\dd (l\cos\theta)}{\dd t} &= \dot l\cos\theta + l\,\frac{\dd\cos\theta}{\dd\theta}\,\frac{\dd\theta}{\dd t} = -l\sin\theta\,\dot\theta .
\end{align}
Hence
\begin{equation}
  \vb = \begin{bmatrix} \dot x + l\cos\theta\,\dot\theta \\ -l\sin\theta\,\dot\theta \end{bmatrix}.
  \label{eq:vb}
\end{equation}

\subsection{Kinetic energy}

The kinetic energy of a rigid body is the sum of the translational energy of
its centre of mass and the rotational energy about its centre of mass,
$\frac12 m v^2 + \frac12 I\omega^2$. For the two bodies this gives four terms.

\paragraph{Wheel, translation.}
The centre of the wheel is the axle, at $(x,0)$, so its velocity is $\dot x$:
\begin{equation}
  T_{\mathrm{w,t}} = \frac12\mw\dot x^2 .
\end{equation}

\paragraph{Wheel, rotation.}
The angular velocity of the wheel is $\dot\varphi$, which by
\cref{eq:rolling-vel} equals $\dot x/r$:
\begin{equation}
  T_{\mathrm{w,r}} = \frac12\Iw\dot\varphi^2 = \frac12\Iw\left(\frac{\dot x}{r}\right)^2 .
\end{equation}

\paragraph{Body, translation.}
The squared speed is the sum of the squares of the two components of
\cref{eq:vb}:
\begin{align}
  v_{\mathrm b}^2 &= (\dot x + l\cos\theta\,\dot\theta)^2 + (-l\sin\theta\,\dot\theta)^2 \nonumber\\
  &= \dot x^2 + 2l\cos\theta\,\dot x\dot\theta + l^2\cos^2\theta\,\dot\theta^2 + l^2\sin^2\theta\,\dot\theta^2 \nonumber\\
  &= \dot x^2 + 2l\cos\theta\,\dot x\dot\theta + l^2\dot\theta^2\left(\cos^2\theta + \sin^2\theta\right) \nonumber\\
  &= \dot x^2 + 2l\cos\theta\,\dot x\dot\theta + l^2\dot\theta^2 ,
\end{align}
where the last step uses $\cos^2\theta + \sin^2\theta = 1$. Therefore
\begin{equation}
  T_{\mathrm{b,t}} = \frac12\mb v_{\mathrm b}^2
  = \frac12\mb\left(\dot x^2 + 2l\cos\theta\,\dot x\dot\theta + l^2\dot\theta^2\right).
\end{equation}

\paragraph{Body, rotation.}
The orientation of the body is the angle $\theta$, measured from the vertical,
which is a fixed direction; the angular velocity of the body is therefore
$\dot\theta$:
\begin{equation}
  T_{\mathrm{b,r}} = \frac12\Ib\dot\theta^2 .
\end{equation}

\paragraph{Total kinetic energy.}
Summing the four terms,
\begin{equation}
  T = \frac12\mw\dot x^2 + \frac12\Iw\left(\frac{\dot x}{r}\right)^2
      + \frac12\mb\left(\dot x^2 + 2l\cos\theta\,\dot x\dot\theta + l^2\dot\theta^2\right)
      + \frac12\Ib\dot\theta^2 .
\end{equation}
Expanding the body term,
\begin{equation}
  T = \frac12\mw\dot x^2 + \frac12\frac{\Iw}{r^2}\dot x^2 + \frac12\mb\dot x^2
      + \mb l\cos\theta\,\dot x\dot\theta + \frac12\mb l^2\dot\theta^2 + \frac12\Ib\dot\theta^2 ,
\end{equation}
and collecting the terms in $\dot x^2$, $\dot x\dot\theta$ and $\dot\theta^2$:
\begin{equation}
  T = \frac12\dot x^2\left(\mw + \frac{\Iw}{r^2} + \mb\right)
      + \mb l\cos\theta\,\dot x\dot\theta
      + \frac12\dot\theta^2\left(\mb l^2 + \Ib\right).
  \label{eq:T}
\end{equation}
For brevity, the three coefficients are denoted
\begin{equation}
  a = \mw + \frac{\Iw}{r^2} + \mb , \qquad
  b = \mb l\cos\theta , \qquad
  c = \mb l^2 + \Ib ,
  \label{eq:abc}
\end{equation}
so that $T = \frac12 a\dot x^2 + b\,\dot x\dot\theta + \frac12 c\,\dot\theta^2$.

\subsection{Potential energy}

The reference height $z=0$ is the wheel axle. The centre of the wheel stays at
this height on flat ground and does not contribute. The body's centre of mass
is at height $l\cos\theta$, by \cref{eq:pb}:
\begin{equation}
  V = \mb g\,l\cos\theta .
  \label{eq:V}
\end{equation}

\subsection{Generalised forces}

The motor torque enters the equations of motion as a generalised force
$Q_i$ for each coordinate, defined by the virtual work $\delta W$ done by the
motor for virtual displacements $\delta x$ and $\delta\theta$:
\begin{equation}
  \delta W = Q_x\,\delta x + Q_\theta\,\delta\theta .
  \label{eq:dW-def}
\end{equation}
The work of a torque is the torque multiplied by the rotation of the body it
acts on. The motor applies $+\tau$ to the wheel, which rotates by
$\delta\varphi$, and $-\tau$ to the body, which rotates by $\delta\theta$
(\cref{fig:model-act}):
\begin{equation}
  \delta W = \tau\,\delta\varphi - \tau\,\delta\theta = \tau\,\delta(\varphi - \theta),
\end{equation}
i.e.\ the work depends on the rotation of the wheel relative to the body. The
virtual rotation of the wheel follows from the rolling constraint
\cref{eq:rolling}, with $\delta r = 0$:
\begin{equation}
  \delta x = \delta(\varphi r) = r\,\delta\varphi + \varphi\,\delta r = r\,\delta\varphi
  \qquad\Longrightarrow\qquad
  \delta\varphi = \frac{\delta x}{r}.
\end{equation}
Substituting,
\begin{equation}
  \delta W = \frac{\tau}{r}\,\delta x - \tau\,\delta\theta ,
\end{equation}
and comparing with \cref{eq:dW-def}:
\begin{equation}
  Q_x = \frac{\tau}{r}, \qquad Q_\theta = -\tau .
  \label{eq:Q}
\end{equation}
$Q_x$ is a force (\si{N}) and $Q_\theta$ a torque (\si{N.m}), consistently with
$x$ being a length and $\theta$ an angle. The contact forces between wheel and
ground do no work, because the contact point does not slide; gravity is
already included in $V$.

\subsection{Euler--Lagrange equations}

The Lagrangian is $L = T - V$. From \cref{eq:T,eq:V},
\begin{equation}
  L = \frac12 a\dot x^2 + \mb l\cos\theta\,\dot x\dot\theta + \frac12 c\,\dot\theta^2 - \mb g l\cos\theta ,
  \label{eq:L}
\end{equation}
and for each coordinate $q_i\in\{x,\theta\}$ the Euler--Lagrange equation is
\begin{equation}
  \frac{\dd}{\dd t}\left(\frac{\partial L}{\partial\dot q_i}\right) - \frac{\partial L}{\partial q_i} = Q_i .
  \label{eq:EL}
\end{equation}

\paragraph{Partial derivatives.}
Each partial derivative is taken with respect to one variable, keeping all the
others constant. With respect to the velocities:
\begin{align}
  \frac{\partial L}{\partial\dot x} &= a\dot x + \mb l\cos\theta\,\dot\theta , \label{eq:dLdxd}\\
  \frac{\partial L}{\partial\dot\theta} &= \mb l\cos\theta\,\dot x + c\,\dot\theta . \label{eq:dLdthd}
\end{align}
With respect to the positions: $L$ does not contain $x$, and $\theta$ appears
in the coupling term and in the potential energy, through $\cos\theta$:
\begin{align}
  \frac{\partial L}{\partial x} &= 0 , \\
  \frac{\partial L}{\partial\theta} &= -\mb l\sin\theta\,\dot x\dot\theta + \mb g l\sin\theta .
  \label{eq:dLdth}
\end{align}

\paragraph{Time derivatives.}
The total time derivative follows the motion, so every quantity that depends
on time ($x$, $\theta$, $\dot x$, $\dot\theta$) is differentiated; only the
parameters are constant. In \cref{eq:dLdxd}, $a$ is constant, so
$\frac{\dd}{\dd t}(a\dot x) = a\ddot x$; the second term is the product of the
two functions of time $\cos\theta$ and $\dot\theta$, and
$\frac{\dd}{\dd t}\cos\theta = -\sin\theta\,\dot\theta$:
\begin{align}
  \frac{\dd}{\dd t}\frac{\partial L}{\partial\dot x}
  &= a\ddot x + \mb l\left(-\sin\theta\,\dot\theta\cdot\dot\theta + \cos\theta\,\ddot\theta\right) \nonumber\\
  &= a\ddot x - \mb l\sin\theta\,\dot\theta^2 + \mb l\cos\theta\,\ddot\theta .
\end{align}
In the same way, from \cref{eq:dLdthd}:
\begin{align}
  \frac{\dd}{\dd t}\frac{\partial L}{\partial\dot\theta}
  &= \mb l\left(-\sin\theta\,\dot\theta\,\dot x + \cos\theta\,\ddot x\right) + c\,\ddot\theta \nonumber\\
  &= -\mb l\sin\theta\,\dot\theta\dot x + \mb l\cos\theta\,\ddot x + c\,\ddot\theta .
\end{align}

\paragraph{Equation of motion for $x$.}
Substituting into \cref{eq:EL} with $Q_x = \tau/r$:
\begin{equation}
  a\ddot x - \mb l\sin\theta\,\dot\theta^2 + \mb l\cos\theta\,\ddot\theta - 0 = \frac{\tau}{r}.
\end{equation}

\paragraph{Equation of motion for $\theta$.}
Substituting into \cref{eq:EL} with $Q_\theta = -\tau$, the minus sign in front
of $\partial L/\partial\theta$ applies to both terms of \cref{eq:dLdth}:
\begin{align}
  -\mb l\sin\theta\,\dot\theta\dot x + \mb l\cos\theta\,\ddot x + c\,\ddot\theta
  - \left(-\mb l\sin\theta\,\dot x\dot\theta + \mb g l\sin\theta\right) &= -\tau \nonumber\\
  -\mb l\sin\theta\,\dot\theta\dot x + \mb l\cos\theta\,\ddot x + c\,\ddot\theta
  + \mb l\sin\theta\,\dot x\dot\theta - \mb g l\sin\theta &= -\tau .
\end{align}
The first and fourth terms cancel, leaving
\begin{equation}
  \mb l\cos\theta\,\ddot x + c\,\ddot\theta - \mb g l\sin\theta = -\tau .
\end{equation}

\paragraph{Summary.}
With $a$ and $c$ from \cref{eq:abc}, the equations of motion are
\begin{align}
  a\ddot x + \mb l\cos\theta\,\ddot\theta - \mb l\sin\theta\,\dot\theta^2 &= \frac{\tau}{r} , \label{eq:eom-x}\\
  \mb l\cos\theta\,\ddot x + c\,\ddot\theta - \mb g l\sin\theta &= -\tau . \label{eq:eom-th}
\end{align}

\subsection{Matrix form}

\Cref{eq:eom-x,eq:eom-th} are written in the form
\begin{equation}
  \M(\q)\,\qdd + \h(\q,\qd) = \B\,\tau ,
  \label{eq:manipulator}
\end{equation}
where row 1 corresponds to the equation of $x$ and row 2 to the equation of
$\theta$. In each equation, the coefficients of the accelerations $\ddot x$ and
$\ddot\theta$ form the corresponding row of the mass matrix $\M$; the remaining
terms on the left-hand side form the vector $\h$; the coefficients of $\tau$
on the right-hand side form the vector $\B$:
\begin{equation}
  \underbrace{\begin{bmatrix}
     \mw + \dfrac{\Iw}{r^2} + \mb & \mb l\cos\theta \\[8pt]
     \mb l\cos\theta & \mb l^2 + \Ib
  \end{bmatrix}}_{\M(\q)}
  \begin{bmatrix} \ddot x \\ \ddot\theta \end{bmatrix}
  +
  \underbrace{\begin{bmatrix}
     -\mb l\sin\theta\,\dot\theta^2 \\[4pt]
     -\mb g l\sin\theta
  \end{bmatrix}}_{\h(\q,\qd)}
  =
  \underbrace{\begin{bmatrix} 1/r \\[4pt] -1 \end{bmatrix}}_{\B}\tau .
  \label{eq:result}
\end{equation}

The same mass matrix follows directly from the kinetic energy. Writing
$T = \frac12\qd^{\mathsf T}\M\qd$ for a generic $2\times2$ matrix $\M$,
\begin{equation}
  \frac12\qd^{\mathsf T}\M\qd
  = \frac12\left(M_{11}\dot x^2 + (M_{12}+M_{21})\,\dot x\dot\theta + M_{22}\dot\theta^2\right),
\end{equation}
and comparing with $T = \frac12 a\dot x^2 + b\,\dot x\dot\theta + \frac12 c\,\dot\theta^2$
gives $M_{11} = a$, $M_{22} = c$ and, for a symmetric matrix,
$M_{12} = M_{21} = b$, in agreement with \cref{eq:result}.

\subsection{Properties of the model}

\paragraph{Mass matrix.}
$\M$ is symmetric, $\M^{\mathsf T} = \M$. Its diagonal entries $a$ and $c$
are constant; only the off-diagonal entry $b = \mb l\cos\theta$ depends on the
configuration. Its determinant,
\begin{equation}
  \det\M = ac - b^2
  = \left(\mw + \frac{\Iw}{r^2}\right)\left(\mb l^2 + \Ib\right) + \mb\Ib + \mb^2 l^2\sin^2\theta ,
\end{equation}
is positive for every $\theta$, so $\M$ is always invertible.

\paragraph{Equilibrium.}
For the robot upright and at rest with no input ($\theta = 0$,
$\dot x = \dot\theta = 0$, $\tau = 0$):
\begin{equation}
  h_1 = -\mb l\sin 0\cdot 0^2 = 0, \qquad
  h_2 = -\mb g l\sin 0 = 0, \qquad
  \B\tau = \bm 0 .
\end{equation}
\Cref{eq:manipulator} reduces to $\M\qdd = \bm 0$ and, since $\M$ is
invertible, $\qdd = \bm 0$: the upright configuration at rest is an
equilibrium.

\paragraph{Stability of the equilibrium.}
With the robot at rest, $\tau = 0$ and a small tilt $\theta > 0$,
\cref{eq:manipulator} gives
$\M\qdd = -\h = \begin{bmatrix} 0 & \mb g l\sin\theta\end{bmatrix}^{\mathsf T}$,
hence
\begin{equation}
  \ddot\theta = \frac{a\,\mb g l\sin\theta}{\det\M} > 0 .
\end{equation}
The angular acceleration has the same sign as the tilt, so the tilt grows: the
equilibrium is unstable. Consistently, the potential energy
\cref{eq:V} is maximum at $\theta = 0$.

\paragraph{Underactuation.}
$\B = \begin{bmatrix} 1/r & -1\end{bmatrix}^{\mathsf T}$ is a single non-zero
column, so $\operatorname{rank}\B = 1$, while the system has two degrees of
freedom. The input $\tau$ always acts on both coordinates in the fixed ratio
$Q_x : Q_\theta = 1/r : -1$, and $\ddot x$ and $\ddot\theta$ cannot be
assigned independently: the system is underactuated.
